\documentclass[11pt]{article}
\usepackage[a4paper,margin=1in]{geometry}
\usepackage[T1]{fontenc}
\usepackage[utf8]{inputenc}
\usepackage{times}
\usepackage{parskip}

\begin{document}

\begin{center}
{\Large\textbf{Statement of Novelty}}\\[1em]
\small{Fixed Points and Certified Boundedness of Unrolled Born/Contrast-Source Inverse Scattering Networks}\\[0.5em]
\small{Longgang Wang, Guoqiang Wang, and Nianfeng Shi}
\end{center}

\vspace{1em}

\noindent\textbf{Statement of Novelty.} The novelty lies in an obstruction rather than in a performance claim. For the Born/contrast-source family of unrolled inverse-scattering networks we show that the Jacobian block coupling contrast perturbations into the contrast-source channel equals a scaled diagonal of the total field, so a closed-form lower bound on the layer's largest singular value follows at once. That bound implies a step-size threshold above which no correction confined to the contrast channel can make the layer contractive, which rules out the per-layer Euclidean condition used elsewhere in the unrolling literature. Certification is therefore rebuilt layer by layer: a field-subsystem contraction condition, an invariant compact set under a nonexpansive projection, and an exact finite-depth noise account that keeps the non-normal factor a geometric series omits. The impact for readers is a reusable template that separates what can be certified from what is only measured, and that identifies the projection, not the learned correction, as the ingredient that bounds the rollout.

\end{document}
